\documentclass[11pt]{article}
\usepackage[a4paper,margin=2.4cm]{geometry}
\usepackage{amsmath,amssymb}
\usepackage{booktabs}
\usepackage{graphicx}
\usepackage{microtype}
\usepackage{natbib}
\usepackage[hidelinks]{hyperref}
\usepackage{enumitem}

\newcommand{\methodname}{VietHerb-MIL}

\title{From Plant-Level Traditional Use to Compound-Level Target Activity:\
Leakage-Controlled Multiple-Instance Learning for Vietnamese Medicinal Plants}
\author{Minh Nhat\\
\small Corresponding author: \texttt{minhnhatpro2504@gmail.com}}
\date{21 October 2026}

\begin{document}
\maketitle

\begin{abstract}
\textbf{Background:}
Traditional-medicine databases connect plants to reported uses and plants to
constituent metabolites, whereas modern bioactivity databases measure activity
at the compound--target level. Naively assigning every plant-level use to every
reported metabolite creates false labels and obscures which chemical evidence,
if any, supports the use. We study whether Vietnamese ethnopharmacological
annotations can improve molecular-target ranking without making this
cross-granularity assumption.

\textbf{Methods:}
We specify a confirmatory workflow linking a maintainer-authorized, versioned
VietHerb snapshot to ChEMBL 37 activities for human DPP4,
acetylcholinesterase, and COX-2. Each plant is represented as a bag of
standardized metabolites, with gated-attention multiple-instance learning and
positive--unlabelled treatment of absent use annotations. Before restricted
VietHerb data were available, we validated the activity branch on a hashed live
ChEMBL snapshot using connected-component scaffold splits, five seeds, four
fingerprint baselines, and dependence-aware paired bootstrap intervals.

\textbf{Results:}
The exploratory snapshot contained 23,602 retrieved records, 16,941 records
after assay filtering, 14,051 aggregated compound--target pairs, and 14,029
unique molecules. No molecule, InChIKey, Murcko scaffold, or connected
component crossed train, validation, and test partitions. Extra Trees achieved
the highest macro PR-AUC (0.8842 $\pm$ 0.0010). Shared activity-only learning
achieved 0.8635 $\pm$ 0.0019 versus 0.8546 $\pm$ 0.0045 for separate neural
models; the paired difference was 0.0089 (hierarchical bootstrap 95\% CI,
0.0034--0.0154). These results validate the activity pipeline but do not test
the ethnopharmacological MIL hypothesis.

\textbf{Conclusions:}
The implemented workflow provides a leakage-audited basis for testing whether
plant-level knowledge adds compound-level ranking value. Exploratory ChEMBL
results support the engineering validity of the activity branch, while the
central MIL comparison remains prospectively locked pending authorized
VietHerb data and expert-audited therapeutic-use mappings.
\end{abstract}

\section{Introduction}

Natural products remain an important source of chemical diversity for drug
discovery, but their biological evidence is heterogeneous and frequently
recorded at a different level from the molecule being screened. VietHerb was
created to organize Vietnamese medicinal plants together with metabolites,
therapeutic effects, morphology, and geography. Its original release reported
2,881 species, 10,887 metabolites, and many-to-many species--metabolite and
species--effect relations \citep{nguyenvo2019vietherb}. This local resource is
scientifically valuable, but a reported use of a plant is neither evidence that
every constituent is active nor proof of clinical efficacy.

Prior computational ethnopharmacology already links traditional symptoms,
efficacies, herbs, ingredients, and molecular targets. HTINet and HTINet2 use
heterogeneous networks and knowledge graphs for herb--target prediction
\citep{wang2019htinet,duan2024htinet2}; other systems use clinical symptom
profiles or multimodal ingredient--target information
\citep{cheung2026novelhti,liang2025multiiti}. Consequently, the broad claim that
traditional knowledge has not been used for target prediction is untenable.
The unresolved issue addressed here is narrower: VietHerb annotations are
plant-level labels, while ChEMBL observations are compound-level measurements.
This is naturally a multiple-instance learning (MIL) problem, in which a plant
is a bag and its reported metabolites are instances
\citep{ilse2018attention,doran2012mil}.

Reliable evaluation is equally important. Ligand benchmarks can reward
memorization when close analogues cross partitions \citep{wallach2018memorization},
and conventional scaffold splitting does not eliminate every form of chemical
redundancy. Traditional-medicine data add botanical leakage because the same
metabolite can occur in several species. Missing compound--target and
plant--use edges also cannot automatically be interpreted as negatives, making
positive--unlabelled (PU) learning a relevant sensitivity analysis
\citep{kiryo2017nnpu,wang2026prodti}.

We therefore make four contributions. First, we define an auditable data model
that preserves compound, plant, therapeutic-concept, assay, and provenance
relations. Second, we formulate cross-granularity supervision with attention
MIL instead of label broadcasting. Third, we introduce split manifests that
keep identical structures, related chemistry, and plant components disjoint.
Fourth, we isolate the value of ethnopharmacological information with matched
architectures and falsification controls. The preregistered primary question is
whether plant-bag supervision improves compound-level out-of-distribution
PR-AUC beyond activity-only multi-task learning.

\section{Materials and methods}

\subsection{Data sources and governance}

VietHerb provides plant--metabolite and plant--therapeutic-effect relations
\citep{nguyenvo2019vietherb}. Because no explicit database redistribution
licence was found during our audit, the confirmatory analysis requires a
maintainer-authorized export. Raw rows will not be redistributed without
written permission. Published artifacts will otherwise contain code, schemas,
source identifiers, checksums, aggregate counts, and split-generation logic.
Traditional-use records are treated as hypothesis-generating cultural and
historical evidence, not as validated biomedical claims.

For the confirmatory study, measured activities will be extracted from the
versioned ChEMBL 37 SQLite release. The results reported here use a live REST
snapshot only for exploratory engineering. ChEMBL is a
manually curated, large-scale bioactivity resource with explicit assay and
document provenance \citep{zdrazil2024chembl}. LOTUS may be used only for an
optional occurrence/provenance sensitivity analysis
\citep{rutz2022lotus}. Aggregated natural-product databases are not assumed to
provide independent activity ground truth \citep{chandrasekhar2025coconut}.

\subsection{Target-activity cohort}

The pilot targets are human dipeptidyl peptidase 4 (DPP4, CHEMBL284),
acetylcholinesterase (AChE, CHEMBL220), and prostaglandin G/H synthase 2
(COX-2/PTGS2, CHEMBL230). The primary cohort contains high-confidence,
single-protein biochemical binding assays with standardized IC$_{50}$ or
$K_i$ values in nM. Functional assays form a separate sensitivity cohort.
Records with data-validity warnings, duplicate flags, unresolved structures, or
ambiguous target assignments are excluded. Exact standardized
structure--target replicates are aggregated by median pChEMBL; pairs spanning
at least one log unit are flagged as contradictory and excluded from the
primary analysis.

The prespecified binary endpoint is active at pChEMBL $\geq 6$. Thresholds 5
and 7 and continuous pActivity models are sensitivity analyses. Unmeasured
compound--target pairs remain masked. All stereochemistry and structure
standardization decisions, rejected rows, release metadata, and SHA-256 hashes
are retained in a machine-readable manifest.

\subsection{Therapeutic-use normalization}

Original text is preserved. Candidate mappings are retrieved from frozen
MeSH/Mondo releases, after which a language model may rank candidates but may
not accept a label automatically. The stored record includes source span,
ontology version, candidate list, confidence, model/prompt hash, abstention,
and reviewer decisions. Two pharmacognosy or pharmacy reviewers independently
annotate a prespecified stratified sample before adjudication. We report exact
and ancestor-aware precision, recall, F1, raw agreement, and Cohen's $\kappa$;
sample size is selected for a target confidence-interval width rather than fixed
solely by convenience.

The first model uses three coarse, clinically interpretable families:
metabolic/diabetes, cognition/neurological use, and pain/inflammation. ATC is
not used as the primary disease ontology because it classifies active drug
substances rather than diseases or symptoms.

\subsection{Molecular and botanical split construction}

Molecules are standardized globally and assigned once across all targets and
plants. Two locked regimes answer different questions. The chemical-OOD regime
groups molecules by a prespecified fingerprint clustering rule and reports
train--test nearest-neighbour Tanimoto distributions. A scaffold component split
is retained as a reproducible secondary regime. The botanical cold-start regime
holds out plant taxa and explicitly handles metabolites shared between taxa.

For the strict joint split, molecule nodes are joined when they share a chemical
group or occur in the same plant. Entire transitive connected components are
assigned to train, validation, or test. Thus, if compound A is chemically linked
to B and B occurs with C in one species, all three remain in one partition. No
feature transformation, loss weight, class prior, or hyperparameter is fitted
using test data.

\subsection{Multiple-instance multi-task model}

Let $x_{ij}$ denote metabolite $j$ reported in plant $i$, and let
$h_{ij}=f_\theta(x_{ij})$ be its shared molecular representation. The primary
heads predict measured compound--target activities with masked, class-weighted
binary cross entropy. For each plant bag, gated attention computes
\begin{align}
e_{ij} &= w^\top\left[\tanh(Vh_{ij}) \odot \sigma(Uh_{ij})\right],\\
a_{ij} &= \frac{\exp(e_{ij})}{\sum_k \exp(e_{ik})},\\
z_i &= \sum_j a_{ij}h_{ij},
\end{align}
followed by one plant-use head per therapeutic family
\citep{ilse2018attention}. Attention weights provide instance attribution but
are not interpreted as causal explanations.

An observed plant use is positive and absence is unlabelled. Auxiliary heads
therefore use non-negative PU risk \citep{kiryo2017nnpu}. For class prior
$\pi$, positive risk $R_p^+$, positive-as-negative risk $R_p^-$, and unlabelled
negative risk $R_u^-$, the unbiased estimate is
\begin{equation}
R_{\mathrm{PU}}=\pi R_p^+ + R_u^- - \pi R_p^- ,
\end{equation}
with the non-negative correction applied when the empirical negative component
falls below zero. Priors are frozen from training-only estimation and varied in
sensitivity analyses. The total objective is
\begin{equation}
\mathcal{L}=\mathcal{L}_{\mathrm{activity}}+
\lambda\mathcal{L}_{\mathrm{plant\text{-}MIL}}.
\end{equation}

The prespecified compute-feasible confirmatory encoder is a radius-2, 2,048-bit chiral Morgan
fingerprint followed by a two-layer MLP. This deliberately separates the
scientific contribution from an expensive foundation model. A D-MPNN/Chemprop
baseline \citep{yang2019chemprop} and one pretrained molecular encoder are
secondary comparisons using identical split IDs.

\subsection{Baselines and falsification controls}

Prespecified models are: per-target fingerprint logistic regression; random forest
or gradient boosting; separate neural target models; shared activity-only
multi-task learning; naive plant-label broadcasting; the proposed MIL model;
and optional domain importance weighting. Controls shuffle use labels among
training plants while preserving prevalence, replace uses with equally sized
random auxiliary labels, compare expert-verified with confidence-weighted
mappings, and equalize parameter counts and update counts. The primary
comparison is MIL versus the architecture-matched activity-only model. Naive
broadcasting is an ablation, not the proposed method.

\subsection{Training and statistical analysis}

The MLP uses hidden dimensions 1,024 and 512, dropout 0.25, AdamW with learning
rate $3\times10^{-4}$ and weight decay $10^{-4}$, batch size 256, and early
stopping patience 12 on validation macro PR-AUC. Auxiliary weights
$\{0.1,0.3,1.0\}$ are selected on validation data. Seeds are 13, 37, 71, 101,
and 137.

The confirmatory primary endpoint is the paired change in macro test PR-AUC. We additionally
report target-wise PR-AUC and prevalence, BEDROC ($\alpha=20$), EF1\% and EF5\%,
precision/recall at fixed experimental budgets, AUROC, Brier score, and
calibration. PR-AUC is emphasized for imbalanced screening
\citep{saito2015pr}, while BEDROC measures early recognition
\citep{truchon2007bedroc}. Paired 95\% intervals use 1,000 bootstrap replicates
with structure cluster as the resampling unit for compound endpoints and plant
as the unit for bag endpoints. Target-level secondary tests use
Benjamini--Hochberg correction.

For the exploratory activity-only validation, the live snapshot was retrieved
on 20 October 2026 and bound to the configuration by SHA-256. Logistic
regression and Extra Trees were compared with separate target-specific MLPs and
an architecture-matched shared MLP. Neural and Extra Trees estimates use the
five prespecified seeds; deterministic logistic regression is repeated only to
retain the same reporting schema. All exploratory manifests are marked
\texttt{confirmatory\_eligible=false}. Runs used Python 3.11.9, RDKit
2025.09.5, scikit-learn 1.8.0, and PyTorch 2.5.1 with CUDA 12.1 on an NVIDIA
RTX 4060 Laptop GPU with 8 GB memory.

\section{Results}

\subsection{Cohort and shift audit}

The REST snapshot yielded 23,602 raw target records. Restriction to eligible
binding measurements retained 16,941 rows, which became 14,051
structure--target pairs after replicate aggregation and 14,029 globally unique
molecules after standardization. Connected-component assignment produced
9,820 training, 2,105 validation, and 2,104 test molecules spanning 5,433
Murcko scaffolds. Pairwise overlap between every split was zero for molecule
identifier, InChIKey, scaffold, and connected-component identifier. Table~1
summarizes the target observations. No VietHerb row was used in this analysis.

\begin{table}[ht]
\centering
\caption{Exploratory ChEMBL REST-snapshot cohort after curation. Active denotes
pChEMBL $\geq 6$. This is not the confirmatory ChEMBL 37/VietHerb cohort.}
\begin{tabular}{lrrrr}
\toprule
Endpoint & Measurements & Molecules & Active (\%) & Test positives\\
\midrule
DPP4 & 4,156 & 4,156 & 83.7 & 522\\
AChE & 5,784 & 5,784 & 49.7 & 431\\
COX-2 & 4,111 & 4,111 & 55.8 & 344\\
All targets & 14,051 & 14,029 & 61.5 & 1,297\\
Plant-use bags & 0 & 0 & --- & 0\\
\bottomrule
\end{tabular}
\end{table}

The absence of split overlap satisfies the implemented identity and scaffold
leakage gates. Nearest-neighbour shift, bag-size distributions, plants per
metabolite, and botanical component sizes cannot be estimated until the
authorized VietHerb snapshot is frozen; these remain confirmatory diagnostics
rather than unreported exploratory outcomes.

\subsection{Exploratory activity-only comparison}

Extra Trees ranked first on macro PR-AUC, BEDROC20, AUROC, and Brier score
(Table~2). The shared MLP exceeded the separate neural models by 0.0089 macro
PR-AUC. A hierarchical bootstrap resampling both seed and scaffold component
gave a 95\% interval of 0.0034--0.0154. Mean target-wise differences for shared
minus separate MLPs were 0.0102 for DPP4, 0.0027 for AChE, and 0.0139 for
COX-2; their across-seed standard deviations were 0.0056, 0.0068, and 0.0134,
respectively. This comparison measures activity-only parameter sharing and is
not evidence for the proposed auxiliary plant supervision.

\begin{table}[ht]
\centering
\small
\caption{Exploratory scaffold-component test performance, reported as mean
$\pm$ standard deviation across five seeds. EF1\% is an enrichment factor, not
a percentage.}
\begin{tabular}{lcccc}
\toprule
Model & Macro PR-AUC & BEDROC20 & EF1\% & Brier\\
\midrule
Logistic regression & $0.8358\pm0.0000$ & $0.9013\pm0.0000$ & $1.5992\pm0.0000$ & $0.1902\pm0.0000$\\
Extra Trees & $0.8842\pm0.0010$ & $0.9532\pm0.0031$ & $1.6434\pm0.0335$ & $0.1505\pm0.0005$\\
Separate target MLPs & $0.8546\pm0.0045$ & $0.9315\pm0.0128$ & $1.6584\pm0.0000$ & $0.1791\pm0.0037$\\
Shared activity MLP & $0.8635\pm0.0019$ & $0.9460\pm0.0063$ & $1.6584\pm0.0000$ & $0.1777\pm0.0031$\\
\bottomrule
\end{tabular}
\end{table}

\subsection{Implementation and confirmatory readiness}

All 32 unit tests passed, including masked-loss, nnPU-risk, attention-MIL,
metric, and connected-component leakage invariants. A CUDA smoke run exercised
the activity-only, naive-broadcast, MIL, and shuffled-MIL paths end to end on
deterministic synthetic bags. Its manifest is explicitly marked
\texttt{scientific\_results=false}; therefore synthetic performance is not
reported. Botanical cold starts, PU-prior and auxiliary-weight sensitivity,
and attention stability require real bags and remain locked confirmatory
analyses.

\section{Discussion}

The exploratory results establish that the curation, splitting, training, and
evaluation stack works on measured bioactivity. They also show that model
selection remains consequential: Extra Trees outperformed both neural
alternatives, while sharing the neural encoder gave a small, consistently
positive macro PR-AUC change relative to separate target models. The confidence
interval quantifies this snapshot-specific comparison; it does not remove bias
from a live database snapshot or elevate an exploratory analysis to a
confirmatory result.

The future principal result must be interpreted as incremental predictive value
under the locked evaluation, not validation of traditional medicine. If MIL improves
over activity-only learning but not over shuffled labels, the gain is consistent
with generic regularization rather than semantic ethnopharmacological signal. If
naive broadcasting performs similarly to MIL, chemical or botanical leakage and
label granularity require further investigation. If all ethnomedical variants
underperform, negative transfer is an informative result and should not be
hidden by selecting favorable targets.

Docking or short molecular dynamics may provide a later structure-based
plausibility assessment, but neither constitutes biological validation.
Docking scores have well-known limitations for affinity ranking
\citep{warren2006docking}. Prospective biochemical assays with positive,
negative, interference, and purity controls are required to claim target
inhibition.

Limitations include incomplete or changing plant databases, uncertain
traditional-use normalization, incomplete chemical inventories for each plant,
assay heterogeneity, class-prior uncertainty in PU learning, and the possibility
that attention is not a faithful explanation. Written data permission and a
domain collaborator are release gates rather than administrative afterthoughts.

\section{Conclusion}

We present a reproducible, leakage-controlled formulation for learning from
plant-level traditional-use labels and compound-level target measurements
without broadcasting each use label to every constituent. A real-data
activity-only benchmark of 14,029 compounds passed all implemented split audits
and showed a modest benefit from shared neural learning, although Extra Trees
remained the strongest exploratory baseline. The central claim---that
ethnopharmacological bag supervision improves out-of-distribution molecular
prioritization---is not yet supported or refuted. It will be tested only after
an authorized VietHerb export, expert-audited mappings, and the version-pinned
ChEMBL 37 cohort pass the prespecified release gates.

\section*{Data and code availability}

Code, configurations, schemas, tests, source locks, and derived manifests are
available at \url{https://github.com/MinhNhat-2504/vietherb-mtl}. ChEMBL-derived
artifacts follow CC BY-SA 3.0. VietHerb row-level data will be released only if
the maintainers provide explicit permission; otherwise the repository will
publish acquisition instructions, hashes, identifiers, and reproducible
transformation code without redistributing restricted rows.

\section*{Author contributions, funding, and competing interests}

Minh Nhat: conceptualization, methodology, software, investigation, data
curation, validation, visualization, and writing. No external funding was used
for the reported exploratory computation. The author declares no competing
interests. Contributions and declarations must be updated if domain reviewers
or experimental collaborators join the confirmatory study.

\section*{Ethics statement}

The reported work used public chemical bioactivity records and deterministic
synthetic bags for software validation; it involved no human participants,
animals, or clinical intervention. Traditional-use annotations are not treated
as clinical evidence.

\section*{Generative-AI assistance disclosure}

A generative-AI coding assistant was used for implementation support, test
generation, literature-search assistance, and language editing. The author
reviewed the source code, citations, numerical outputs, and manuscript claims
and remains responsible for the final work.

\bibliographystyle{plainnat}
\bibliography{references}

\end{document}

